\clearpage
\section{Proof of Theorem \ref{thm:Segala}}
\label{app:ProofSegala}

In our setting, the most interesting implication of Theorem \ref{thm:Segala} is soundness.
We will see that the definition of $\sqsubseteq^t$ has some limitations, notably with respect to
liveness. To address this, we introduce a stronger notion of trace distribution and strengthen
probabilistic forward simulation accordingly. However, the soundness proof of Segala
\cite{Segala1995Compositional} works at the level of $\efrak_{\mu,s}$, not $s$. This is
problematic because our strengthening consists in constraining schedulers; thus, we need to
adapt the proof. If $\scal$ is a probabilistic forward
simulation then it induces a choice function
$\ccal : \scal \times T \rightarrow Sch \times (\scal \rightarrow [0,1])$
which maps any $(q,\mu') \in \scal$ and $(q,l,\mu) \in T$ to the scheduler that induces
$\mu' \xRightarrow{l} \mu''$
and the witness $w$ of $(\mu,\mu'') \in \dcal(\scal)$.

\begin{algorithm}
    \caption{Simulation Algorithm for Schedulers}
    \label{alg:simulation}
    \begin{algorithmic}
        \State \textbf{Input:} $\pcal, \pcal'$ two PLTSs, $\scal$ a probabilistic forward
        simulation of $\pcal$ by $\pcal'$ and $s \in Sch$
        \medskip
        \State $\texttt{e}, \texttt{e'}, \texttt{mu'} \gets q_\star, q'_\star, \delta_{q_\star'}$
        \State $\texttt{acc} \gets \texttt{e'}$
        \State $\texttt{t} \gets s(\texttt{e})$
        \While{$\texttt{t} \ne \bot$}
            \State $\texttt{s'}, \texttt{w} \gets \ccal(last(e),\mu',t)$
            \State $\texttt{t'} \gets \texttt{s'}(\texttt{e'})$
            \While{$\texttt{t'} \ne \bot$}
                \State $\texttt{e'} \gets \texttt{e'}.lbl(\texttt{t'}).out(\texttt{t'})$
                \State $\texttt{acc} \gets \texttt{acc}.lbl(\texttt{t'}).out(\texttt{t'})$
                \State $\texttt{t'} \gets \texttt{s'}(\texttt{e'})$
            \EndWhile
            \State $\texttt{e'}, \texttt{mu'} \gets last(\texttt{e'}),
            \qfrak_{\texttt{mu'},\texttt{s'}}$
            \State $\texttt{e} \gets \texttt{e}.lbl(\texttt{t}).q$ with probability
            $\frac{\sum_{\mu \in \dcal(Q')}w(q,\mu)\cdot\mu(\texttt{e'})}
            {\texttt{mu'}(\texttt{e'})}$
            \State $\texttt{t} \gets \texttt{s}(\texttt{e})$
        \EndWhile
    \end{algorithmic}
\end{algorithm}

Once the choice function $\ccal$ is fixed, Algorithm \ref{alg:simulation} is a sequential
probabilistic program; thus, it can be modelled as a deterministic PLTS. Of course, the
validity of the algorithm's individual instructions should be established rigorously, but
this is not our concern here. The variable $\texttt{e}$ encodes an execution of $\pcal$ and
$\texttt{acc}$ an execution of $\pcal'$ that simulates it. In order to extract a scheduler of
$\pcal'$ from Algorithm \ref{alg:simulation}, we need to identify the instructions deciding which
transitions to execute in $\pcal'$. The most obvious one is
$\texttt{t'} \gets \texttt{s'}(\texttt{e'})$
when $\texttt{t'} \ne \bot$; then there is $\texttt{t} \gets \texttt{s}(\texttt{e})$ when
$\texttt{t} = \bot$. The case where $\texttt{t'} = \bot$ has to be discarded
because another transition will be picked without updating $\texttt{acc}$. When $\texttt{t} = \bot$,
the algorithm stops, so the scheduler on $\pcal'$ also returns $\bot$. After each update
of $\texttt{acc}$, one of the three following events occurs: (i) $\texttt{t}$ is set to $\bot$;
(ii) $\texttt{t'}$ is set to a non-$\bot$ value; or (iii) $\texttt{t'}$ is set to $\bot$
indefinitely.
Cases (i) and (iii) correspond to scheduling $\bot$ and case (ii) corresponds to scheduling
$\texttt{t'}$. Thus, the key values to compute are $reach(e')$, the probability of reaching a
configuration where $\texttt{acc} = e'$, and $sch(e',t)$, the probability of reaching a
configuration where $\texttt{acc} = e$ and $\texttt{t'} = t$. Knowing these, one can define a
scheduler on $\pcal'$ as
\[
    s' : e' \mapsto
        \begin{cases}
            t' \mapsto \frac{sch(e',t')}{reach(e')} \\
            \bot \mapsto 1 - \frac{\sum_{t' \in T'}sch(e',t')}{reach(e')}
        \end{cases}
\]
We conjecture that $\tfrak_{q_\star,s} = \tfrak_{q_\star',s'}$.
